\section{Positive analytic compression of the canonical source}
\label{sec:analytic-source}
This section integrates Note~58 \cite{n58}. It changes the source
representation in Theorem~\ref{thm:all-time}, while retaining its
field table, physical regulator, fixed filter and four-jet consumer.
First fix the upper time $b$ selected there, then a dyadic
$0<\varepsilon\le1$ satisfying every entry of
\eqref{eq:sc-all-source}. A rational upper bound $s\ge\max(0,b)$
or certified recipe integers are inputs to the algorithm. In particular,
fixed-$b$ complexity below excludes the cost of obtaining the upstream
contour and transfer constants. The source remains fixed at all queries.

\subsection{A buffered complex square root of the complete source}
\begin{proposition}[Canonical analytic branch]\label{prop:ac-root}
On $\mathcal D=\{w:\Re w>-1/16,\ |\Im w|<1/16\}$, the
complete source has a holomorphic square root $h$ positive on $[0,\infty)$,
with $|h|<9$ and $|h'(u)|\le288$ for real $u\ge0$.
\end{proposition}
\begin{proof}
Set $E=e^{4w}$ and factor the whole canonical series as
\begin{align}
 \Phi(w)&=A(w)(1+\mathcal R(w)),\notag\\
 A(w)&=\pi e^{5w}(2\pi E-3)e^{-\pi E},\notag\\
 \mathcal R(w)&=\sum_{r\ge2}
 r^2\frac{2\pi r^2E-3}{2\pi E-3}e^{-\pi(r^2-1)E}.
 \label{eq:ac-factor}
\end{align}
Here $|E|>3/4$, $\Re E>93/128>2/3$, using
$e^{-1/4}>3/4$ and $\cos(1/4)\ge31/32$. Since $3<\pi<4$,
$|2\pi E-3|>2|E|$ and $\Re(2\pi E-3)>87/64>0$.
Each relative summand is bounded by
$t_r=6r^4e^{-2(r^2-1)}$. Now $t_2=96e^{-6}$ and
$t_{r+1}/t_r\le(3/2)^4e^{-10}<1/4$, so normal convergence gives
\begin{equation}
 |\mathcal R|\le128e^{-6}<1/2.
 \label{eq:ac-relative-tail}
\end{equation}
Both logarithm factors in
\begin{equation}
 h(w)=\exp\left\{\frac{\log\pi+5w+\Log(2\pi E-3)
                     -\pi E+\Log(1+\mathcal R(w))}{2}\right\}
 \label{eq:ac-root}
\end{equation}
therefore use their principal right-half-plane branches. This proves
$h^2=\Phi$ and real positivity, with no inference of a complex branch
from real positivity alone. If $y=e^{4\Re w}$, then
\[
 |A|\le(32y^{9/4}+12y^{5/4})e^{-93y/32}\le44.
\]
For either exponent $a=9/4,5/4$, the maximum is
$(32a/(93e))^a<1$. Hence $|\Phi|<66$, $|h|<9$.
Cauchy's circle of radius $1/32$ about each $u\ge0$ lies fully
inside $\mathcal D$ and gives $|h'(u)|\le288<512$.
\end{proof}

\subsection{Interpolation, finite coefficient construction and positivity}
Put
\begin{align}
 Q_s&=\log(32/\varepsilon),&
 U&=\lceil4+\log(2+s+Q_s)\rceil,\notag\\
 F_s&=U^4e^{sU^2+4U},&\delta&=\varepsilon/(160F_s),&
 \rho&=1+1/(8U).
 \label{eq:ac-cutoff}
\end{align}
Certified upper integers with bounded slack may replace transcendental
ceilings; if $U$ increases, recompute all later parameters. The full
source-tail proof of Section~\ref{sec:source-conditioning} gives
\begin{equation}
 \int_U^\infty u^je^{su^2+4u}\Phi(u)\,du\le\varepsilon/8,
 \qquad 0\le j\le3.
 \label{eq:ac-full-tail}
\end{equation}
Indeed $s,Q_s\le e^{U-4}$,
$4\pi e^{4U}-2sU-13-j/U\ge\pi e^{4U}$ and
$sU^2+13U-\pi e^{4U}\le-2e^{4U}\le-Q_s$; the complete
tail prefactor is at most one. This is the entire source tail.

Under $w=U(1+v)/2$, the closed Bernstein ellipse of parameter $\rho$
has imaginary extent $U(\rho-\rho^{-1})/4<1/16$ and left
extension $U(\rho+\rho^{-1}-2)/4<1/(256U)$. It lies in
$\mathcal D$. For $g(v)=h(U(1+v)/2)$, the Laurent expansion in
$v=(\zeta+\zeta^{-1})/2$ gives Chebyshev coefficients
$|a_k|\le18\rho^{-k}$. At the roots of $T_{n+1}$ every $T_k$
aliases to $0$ or $\pm T_r$, $r\le n$; thus the interpolation
operator has norm at most one on each individual Chebyshev mode.
Applying this to the convergent tail proves
\[
 \|g-I_ng\|_\infty\le36\rho^{-n}/(\rho-1)
                         =288U\rho^{-n}.
\]
Since $\log\rho\ge1/(8U+1)$, the finite choice
\begin{equation}
 n=\lceil16U\log(1152U/\delta)\rceil,\qquad m=n+1
 \label{eq:ac-degree}
\end{equation}
pays interpolation error at most $\delta/4$.

Here is a finite coefficient recipe, including source arithmetic.
Define
\begin{align}
 \eta_n&=\delta/(16m),&\zeta_n&=\delta/(160m),&
 \beta&=\lceil\log_2(8m/\delta)\rceil,\notag\\
 \alpha_n&=\min\{\eta_n/(4096U),\zeta_n/m^2,1/4\},&
 B_c&=\lceil\log_2(8/\alpha_n)\rceil,&D_{\cos}&=B_c+46.
 \label{eq:ac-node-recipe}
\end{align}
At $\theta_l=(2l+1)\pi/(2m)$ use finite Machin series for $\pi$
to error $2^{-B_c-10}$ and cosine Taylor degree $D_{\cos}$ at
arguments of modulus at most four. Its factorial remainder is at most
$2^{-B_c-10}$. Round to the $B_c$-bit dyadic grid and project onto
$[-1,1]$, obtaining $|\widehat c_l-\cos\theta_l|\le\alpha_n$.
Set $\widehat u_l=U(1+\widehat c_l)/2$ exactly. The branch bound gives
nodal displacement error $\le\eta_n/16$ in $h$. The recurrence for
$T_r$ gives exact rational weights $T_r(\widehat c_l)$; since
$|T_r'|\le r^2$ on $[-1,1]$, weight error is at most $\zeta_n$.
The derivative inequality follows directly from
$\sin(r\theta)/\sin\theta$ having modulus at most $r$.

For each source value choose
\begin{align}
 n_0&=\max\{3,\lceil\sqrt{\log(2^{18}/\eta_n^2)/2}\rceil\},
 &N&=n_0-1,\notag\\
 B_\Phi&=\lceil\log_2(4096/\eta_n^2)\rceil,&
 B_h&=\lceil\log_2(128/\eta_n)\rceil,\notag\\
 T&=\lceil10N^2e^{4U}+13U+1\rceil,&
 D_{\exp}&=\lceil P+8T+4\rceil,\notag\\
 P&=B_\Phi+20+\lceil9\log_2(N+1)+26U/\log2\rceil.
 \label{eq:ac-source-arithmetic}
\end{align}
The complete omitted-value envelope is
$64n_0^6e^{-3n_0^2}\le64e^{-2n_0^2}\le\eta_n^2/4096$.
Note~44's finite Machin/Taylor recipe \cite{n44} with these parameters
computes $\Phi_N(\widehat u_l)$ to error at most
$2^{-B_\Phi}\le\eta_n^2/4096$. To make the arithmetic explicit,
approximate $\pi$ and positive exponentials first; their exact rational
product is the argument of the negative exponential, whose argument
derivative is at most one. The full finite-product error is bounded by
$2^{16}N^9e^{26U}2^{-P}$ before $B_\Phi$-bit rounding.
Degree $D_{\exp}$ pays the largest argument $T$ by the factorial
Taylor bound. Thus no nodal-value oracle is used.

Project that value onto $[0,\infty)$ and perform $B_h$ dyadic
square-root bisections on $[0,16]$. The true value is below $82$;
projection cannot increase its error. Since
$|\sqrt a-\sqrt b|\le\sqrt{|a-b|}$ for $a,b\ge0$, the
source/mode square-root error is at most $\eta_n/\sqrt{2048}$,
and bisection costs at most $\eta_n/8$. Together with node
displacement this gives dyadic $0\le\widehat h_l\le10$ with
$|\widehat h_l-h(U(1+\cos\theta_l)/2)|\le\eta_n$.
The squared accuracy in \eqref{eq:ac-source-arithmetic} is essential;
no inverse lower bound for $\Phi$ near $U$ is required.

Form, by exact rational sums,
\[
 \widehat a_0=m^{-1}\sum_l\widehat h_l,\qquad
 \widehat a_r=(2/m)\sum_l\widehat h_lT_r(\widehat c_l),\quad1\le r<m.
\]
Round each once to grid $2^{-\beta}$, obtaining $\bar a_r$, and set
\begin{equation}
 p(u)=\sum_{r=0}^n\bar a_rT_r(2u/U-1),\qquad
 f(u)=p(u)^2\mathbf1_{[0,U]}(u).
 \label{eq:ac-positive-source}
\end{equation}
The stored polynomial is real, and source positivity is exact.

\begin{theorem}[Compressed source and complete four-jet bound]
\label{thm:ac-source}
The finite source \eqref{eq:ac-positive-source} has exact-heat field
$G(t,z)=\int_0^Ue^{tu^2}f(u)\cos(zu)\,du$ satisfying
\begin{equation}
 |\partial_z^j(G-H)|\le51\varepsilon/256,
 \quad 0\le j\le3,\quad t\le b,\quad x\in\R,\quad |y|\le4.
 \label{eq:ac-four-jets}
\end{equation}
Its source and cutoff are independent of the queried $t,x,y$.
\end{theorem}
\begin{proof}
Data perturbation costs at most $2m\eta_n=\delta/8$ in polynomial
norm; weight perturbation costs $20m\zeta_n=\delta/8$; coefficient
rounding costs $m2^{-\beta}\le\delta/8$. Hence
$\|p-h\|\le5\delta/8<\delta$, and
$|p^2-\Phi|\le(18+\delta)|p-h|\le19|p-h|$.
Each jet's weight is bounded by $u^je^{4u}$, with integral at most
$F_s$. The separate debits are
\begin{center}\small
\begin{tabular}{@{}lr@{}}\toprule
Full analytic interpolation & $19\varepsilon/640$\\
Nodal data (mode, arithmetic, square root and node) & $19\varepsilon/1280$\\
Dyadic-node DCT weights & $19\varepsilon/1280$\\
Stored coefficient rounding & $19\varepsilon/1280$\\
Complete source tail & $\varepsilon/8$\\\bottomrule
\end{tabular}
\end{center}
Their sum is $51\varepsilon/256$. Exact conversion, squaring and
positivity add no approximation error; the factor $19$ already pays
the squaring cross terms. Compact integration permits all
differentiations and gives $G_t=-G_{zz}$ exactly, even if $f$ jumps
at $U$. The source approximation asserts only these four jets;
later evaluation of higher jets does not strengthen this comparison.
\end{proof}

\subsection{Descriptor and construction costs}
\begin{theorem}[Explicit source costs]\label{thm:ac-construction}
With the certified recipe parameters above, a sufficient compressed
descriptor and a normalized-power descriptor satisfy
\begin{align}
 S_{\rm Cheb}&\le64m[\beta+\lceil\log_2(U+m+2)\rceil+1],\notag\\
 L_d&=2\beta+6m+\lceil\log_2(m+1)\rceil+10,\notag\\
 S_{\rm power}&\le64(2m+1)[L_d+\lceil\log_2(U+m+2)\rceil+1].
 \label{eq:ac-descriptors}
\end{align}
Put $K_{\rm int}=2+U+m+N+\beta+B_\Phi+B_c+B_h+P+D_{\exp}+D_{\cos}$.
The source construction needs at most $2^{20}K_{\rm int}^4$
integer/rational operations with operands at most
$2^{20}K_{\rm int}^4$ bits; a sufficient schoolbook bit bound is
$2^{100}K_{\rm int}^{12}$. At fixed $s$, as $q_s=\log(1/\varepsilon)
\to\infty$,
\begin{align}
 n&=O_s(q_s\log q_s),&\beta&=O_s(q_s),\notag\\
 S_{\rm Cheb}&=O_s(q_s^2\log q_s),&
 S_{\rm power}&=O_s(q_s^2\log^2q_s),
 \label{eq:ac-asymptotic-cost}
\end{align}
and the stated constructor has the deliberately coarse bit bound
$O_s(q_s^{60})$.
\end{theorem}
\begin{proof}
The rounded DCT coefficients have modulus at most $21$. For
$Q_r(v)=T_r(2v-1)$, the recurrence gives
$\|Q_r\|_{\ell^1}\le7^r$, so the expanded $p$ has coefficient
norm $<4\cdot7^m$. Squaring gives denominator $2^{2\beta}$ and
numerator lengths at most $L_d$; its exponential coefficient growth
is explicitly included in $6m$. Conversion and squaring use
$O(m^2)$ exact operations. Literal powers of $u$ would additionally
need up to $2n\log_2U$ denominator bits per coefficient.

To control rational operands, evaluate a Taylor polynomial at $a=p/q$
using common denominator $q^D D!$, then round each elementary output
to its assigned dyadic grid before source products. All Taylor arguments
have at most $32K_{\rm int}$ bits; degrees are at most $K_{\rm int}$.
The denominator has at most
$32K_{\rm int}D+D\log_2D+1$ bits and scaled sums at most
$64K_{\rm int}^2$. Machin staging fits $128K_{\rm int}^2$ bits.
DCT denominators can share $m2^{rB_c+B_h}$; their operands,
source products and conversion fit the same enlarged cap. This avoids
unbounded denominator accumulation or an unspecified reduction oracle.
A sufficient operation count is
\[
 2^{16}\{m(N+1)(D_{\exp}+P+1)^2
       +m(D_{\cos}+B_c+1)^2+m^2+mB_h+1\}
       \le2^{20}K_{\rm int}^4.
\]
Schoolbook arithmetic on the preceding operand cap gives the claimed
bit bound, including divisions and bisections. At fixed $s$,
$U=O(\log q_s)$, $\log F_s=O_s(\log^2q_s)$,
$N=O_s(\sqrt{q_s})$, and all accuracy lengths other than degree are
$O_s(q_s)$; $T,D_{\exp},K_{\rm int}=O_s(q_s^5)$ suffice.
Substitution proves the asymptotic bounds.
\end{proof}
For varying $s$, use the finite recipes: $U\le5+\log(2+s+Q_s)$,
$e^{4U}\le e^{20}(2+s+Q_s)^4$ and
$\log F_s=sU^2+4U+4\log U$. No polynomial in $\log b$ or joint
global accuracy/time theorem follows. For the specified fixed-$b$
policy \eqref{eq:sc-radius-cost}, $q_s=\pi R/8+O(\log^2R)$,
so the new sufficient compressed descriptor is $O(R^2\log R)$ bits.
This improves the earlier sufficient clipped-cubic upper bound; it
does not supply a lower bound or a useful finite crossover.

\subsection{Finite Gaussian moments and the cancellation debit}
\begin{proposition}[An entire moment representation]\label{prop:ac-moments}
Write the exact converted source as
\[ f(u)=\sum_{k=0}^{2n}d_k(u/U)^k,\qquad
 M_k(a,c)=\int_0^1v^ke^{av^2+cv}\,dv. \]
Then
\begin{equation}
 G(t,z)=\frac U2\sum_kd_k[M_k(tU^2,izU)+M_k(tU^2,-izU)].
 \label{eq:ac-moments}
\end{equation}
Jets shift $k$ and multiply by powers of $\pm iU$. For $a\ne0$,
$A=\sqrt{-a}$ gives
\begin{align}
 M_0&=\frac{\sqrt\pi}{2A}e^{-c^2/(4a)}
       [\operatorname{erf}(A-c/(2A))-\operatorname{erf}(-c/(2A))],\notag\\
 2aM_1&=e^{a+c}-1-cM_0,\notag\\
 2aM_{k+1}&=e^{a+c}-kM_{k-1}-cM_k\quad(k\ge1).
 \label{eq:ac-moment-recurrence}
\end{align}
At $a=0,c\ne0$, use $M_0=(e^c-1)/c$ and
$M_k=(e^c-kM_{k-1})/c$; at $a=c=0$, use $M_k=1/(k+1)$.
These branches define the same entire field.
\end{proposition}
\begin{proof}
Substitute $u=Uv$, complete the square for $M_0$, and integrate the
derivatives of $e^{av^2+cv}$ and $v^ke^{av^2+cv}$. Changing the
square-root branch changes neither $M_0$ nor the integral. The integral
proves removability and gives the entire expansion
\begin{equation}
 M_k(a,c)=\sum_{r,s\ge0}\frac{a^rc^s}{r!s!(k+2r+s+1)}.
 \label{eq:ac-entire-moments}
\end{equation}
\end{proof}
The symbolic formula has $O(m)$ shared moment nodes, after $O(m^2)$
exact source conversion. This does not count evaluation precision:
near-zero divisions, opposite moment terms and large power coefficients
can cancel. Section~\ref{sec:complete-evaluation} pays those errors
using \eqref{eq:ac-entire-moments} rather than assuming stability of
the error-function or forward recurrence formulas.
